\subsection{Definición de los indicadores espectrales y de la prueba 0--1}
\label{diag:master:definition}

Para una componente postransitoria $u_j$, muestreada uniformemente con
intervalo $\Delta t$, se sustrae su media $\bar u$ y se aplica una ventana
$w_j$. Con $N$ muestras, la transformada discreta es
\begin{equation}
 S_k=\sum_{j=0}^{N-1}w_j(u_j-\bar u)e^{-2\pi\ii jk/N},
 \qquad f_k=\frac{k}{N\Delta t}.
 \label{diag:master:dft}
\end{equation}
Para la ventana de Hann tomamos \(N\geq3\) y
\(w_j=[1-\cos(2\pi j/(N-1))]/2\). La representación de módulo emplea
\(|S_k|/\sum_jw_j\). Para obtener una densidad de potencia partimos de
\(v_j=w_j(u_j-\bar u)\) y definimos la normalización bilateral
\[
 P_k^{(2)}=\frac{\Delta t}{\sum_jw_j^2}|S_k|^2,
 \qquad \Delta f=\frac1{N\Delta t}.
\]
La identidad de Parseval, con la convención de
\cref{diag:master:dft}, es
\(\sum_{k=0}^{N-1}|S_k|^2=N\sum_{j=0}^{N-1}|v_j|^2\). Sustituyéndola,
\begin{align*}
 \sum_{k=0}^{N-1}P_k^{(2)}\Delta f
 &=\frac{\Delta t}{\sum_jw_j^2}\,
   \frac1{N\Delta t}\sum_{k=0}^{N-1}|S_k|^2\\
 &=\frac{\sum_jw_j^2(u_j-\bar u)^2}{\sum_jw_j^2}.
\end{align*}
La integral discreta de la densidad recupera así la media cuadrática
ponderada por la ventana. Para una señal real, las potencias en \(k\) y
\(N-k\) coinciden. Al pasar a una PSD unilateral sumamos cada pareja,
lo que duplica los modos positivos interiores. El modo cero y el de
Nyquist, cuando \(N\) es par, no tienen una pareja distinta y se conservan
sin duplicación. Welch promedia estas estimaciones entre segmentos.
Si \(u\) tiene unidades \(U\), la PSD tiene unidades \(U^2\) por unidad
de frecuencia; el módulo de la transformada conserva unidades \(U\).

Para cuantificar concentración espectral se excluye el modo cero y se
normalizan las potencias $P_k=|S_k|^2$ en los $m$ modos positivos:
\begin{equation}
 p_k=\frac{P_k}{\sum_{\ell=1}^mP_\ell},\qquad
 H=-\frac{\sum_{k=1}^mp_k\log p_k}{\log m}.
 \label{diag:master:entropy}
\end{equation}
Se requiere $m>1$ y potencia total positiva; una señal constante se trata
por separado. Adoptamos \(0\log0=0\). Como \(0<p_k\leq1\) en los modos
con potencia positiva, cada término \(-p_k\log p_k\) es no negativo.
Para la cota superior, sea \(r\leq m\) el número de esos modos. La
concavidad del logaritmo da
\[
 -\sum_{p_k>0}p_k\log p_k
 =\sum_{p_k>0}p_k\log\frac1{p_k}
 \leq\log\!\left(\sum_{p_k>0}p_k\frac1{p_k}\right)
 =\log r\leq\log m.
\]
Dividiendo por \(\log m>0\), resulta \(0\leq H\leq1\). La entropía
es cero si toda la potencia se concentra en un modo y alcanza uno si
\(p_k=1/m\) en los \(m\) modos.
El cociente de potencia dominante es $\max_kp_k$, y la frecuencia
dominante corresponde al modo que lo alcanza. Para comparar ventanas se
conservan paso de muestreo, longitud y ventana; de otro modo cambian la
resolución $1/(N\Delta t)$ y el número de modos usados en $H$.
La deriva relativa usada en el clasificador compara los máximos de
ventanas parciales mediante $(f_{\max}-f_{\min})/f_{\max}$, cuando
$f_{\max}>0$ y todas las frecuencias parciales son positivas. Si alguna
es nula o no puede estimarse, el clasificador asigna deriva infinita y
no acepta el filtro de periodicidad. Una señal transitoria o una resolución insuficiente pueden
producir banda ancha sin demostrar caos.

La prueba 0--1 utiliza al menos $N=100$ muestras finitas de una señal
preparada $\phi_j$, después de
declarar el descarte, submuestreo, eliminación de tendencia y
normalización \cite{gottwald2009}. Para $c\in(0,\pi)$, define
\begin{align}
 p_c(n)&=\sum_{j=1}^n\phi_j\cos(jc),\nonumber\\
 q_c(n)&=\sum_{j=1}^n\phi_j\sin(jc).
 \label{diag:master:translations}
\end{align}
El desplazamiento cuadrático medio a un retardo $n<N$ es
\begin{equation}
 M_c(n)=\frac1{N-n}\sum_{j=1}^{N-n}
 \bigl[\Delta_np_c(j)^2+\Delta_nq_c(j)^2\bigr],
 \label{diag:master:msd}
\end{equation}
donde $\Delta_np_c(j)=p_c(j+n)-p_c(j)$, y análogamente para $q_c$.
Se elimina el término oscilatorio de la media mediante
\begin{equation}
 D_c(n)=M_c(n)-\bar\phi^{\,2}
                 \frac{1-\cos(nc)}{1-\cos c}.
 \label{diag:master:modified-msd}
\end{equation}
Para ver el origen de la corrección, tomamos primero una señal
constante \(\phi_j=\bar\phi\). La suma compleja de las dos traslaciones
es una progresión geométrica:
\[
 \sum_{j=1}^ne^{\ii jc}
 =e^{\ii c}\frac{1-e^{\ii nc}}{1-e^{\ii c}},\qquad
 \left|\sum_{j=1}^ne^{\ii jc}\right|^2
 =\frac{(1-e^{\ii nc})(1-e^{-\ii nc})}
       {(1-e^{\ii c})(1-e^{-\ii c})}
 =\frac{2-2\cos(nc)}{2-2\cos c}.
\]
El factor de fase que aparece al desplazar el intervalo no cambia el módulo.
Como \(c\in(0,\pi)\), el denominador es positivo. Multiplicar la última
expresión por \(\bar\phi^2\) da la contribución constante al desplazamiento;
al restarla obtenemos \cref{diag:master:modified-msd}.

Para el vector de retardos $n=1,\ldots,n_{\max}$ se calcula
\(K_c=\operatorname{corr}(n,D_c(n))\). Si \(L=n_{\max}\),
\(\bar n=L^{-1}\sum_{n=1}^Ln\) y
\(\bar D_c=L^{-1}\sum_{n=1}^LD_c(n)\), la correlación es
\[
 K_c=\frac{\sum_{n=1}^L(n-\bar n)(D_c(n)-\bar D_c)}
 {\sqrt{\sum_{n=1}^L(n-\bar n)^2\,
                \sum_{n=1}^L(D_c(n)-\bar D_c)^2}}.
\]
El valor comunicado es la mediana
de $K_c$ sobre los valores de $c$ elegidos; se conserva esa selección al
comparar ensayos. En el procedimiento de correlación utilizado,
$n_{\max}=\min(\max(20,\lfloor N/10\rfloor),500)$.
Una varianza nula del desplazamiento no admite correlación de Pearson;
se trata como caso degenerado regular. Los umbrales operativos
$K<0.2$, $K>0.8$ y el intervalo intermedio describen, respectivamente,
compatibilidad regular, compatibilidad caótica e indeterminación de ese
diagnóstico finito. Un valor ligeramente negativo es posible por la
correlación muestral. Ruido, transitorios y sobremuestreo afectan el
resultado; la mediana y la comparación de varios intervalos de muestreo
reducen sensibilidades concretas, sin demostrar caos ni ocultedad.

Al comparar trayectorias mantenemos la misma preparación de la señal.
Un cambio de ventana, submuestreo o eliminación de tendencia puede modificar
\(H\) y \(K\) aunque el campo y los parámetros sean los mismos. Por ello
interpretamos estas magnitudes junto con los retornos y los controles de
integración.
